\section{Discussion}



\subsection{Overview of the Findings}



The objective of this study was to examine how TabPFN compares with established gradient-boosted tree algorithms across different training-data regimes. The results show that the answer depends on the dataset, evaluation metric, and amount of available training data.



The experiments do not support a general claim that TabPFN is superior to gradient-boosted tree models. Instead, they demonstrate that TabPFN can provide substantial predictive advantages under some conditions while conventional tree-based models remain highly competitive under others.



The clearest advantage for TabPFN was observed on Bank Marketing. In contrast, CatBoost was the strongest overall model on Adult and Credit-G. This variation is important because it demonstrates that model selection for tabular classification cannot be reduced to choosing a single algorithm based solely on its model family.



The results also reveal an important computational trade-off. TabPFN achieved strong predictive performance in several settings but required considerably more prediction time than the conventional tree-based models in this benchmark.



\subsection{Research Question 1: Does TabPFN Provide an Advantage When Training Data Are Limited?}



The results provide partial support for this question.



TabPFN performs strongly in the lower-data regimes, particularly when compared with XGBoost and LightGBM. This pattern is visible on Adult, Bank Marketing, and Credit-G.



For Adult, TabPFN achieves a mean small-regime ROC-AUC of 0.7993, compared with 0.7304 for XGBoost and 0.6419 for LightGBM. However, CatBoost remains slightly stronger at 0.8015.



For Bank Marketing, the difference is considerably larger. TabPFN achieves a small-regime ROC-AUC of 0.7520, compared with 0.6437 for XGBoost and 0.6195 for LightGBM.



Credit-G provides another example of strong TabPFN performance against XGBoost and LightGBM, although CatBoost remains competitive.



Therefore, the findings support the idea that TabPFN can be particularly competitive when labelled training data are limited. However, the results do not establish that TabPFN is always the best model in low-data settings because CatBoost remains highly competitive, particularly on Adult and Credit-G.



\subsection{Research Question 2: How Does TabPFN's Relative Performance Change as Training Size Increases?}



The results do not support a single universal pattern.



On Adult, TabPFN is highly competitive in smaller and medium training regimes, but CatBoost becomes stronger at larger training sizes. In the large regime, CatBoost achieves a mean ROC-AUC of 0.9182 compared with 0.9116 for TabPFN.



This behaviour is consistent with the possibility that conventional gradient-boosted trees can increasingly exploit additional task-specific training observations.



However, Bank Marketing provides an important counterexample. TabPFN maintains an advantage across the small, medium, and large regimes. Its mean ROC-AUC increases from 0.7520 in the small regime to 0.8907 in the medium regime and 0.9339 in the large regime. In the large regime, this remains higher than the corresponding values for XGBoost, LightGBM, and CatBoost.



Consequently, the results do not support a simple hypothesis that TabPFN's advantage necessarily disappears as more training data become available.



Instead, the effect of training-data size is dataset-dependent.



This is an important finding for the study because it indicates that the usefulness of TabPFN cannot be characterized only by the quantity of available training data. Dataset characteristics also appear to influence whether the model retains an advantage at larger training sizes.



\subsection{Research Question 3: Is the TabPFN Advantage Consistent Across Datasets and Metrics?}



The answer is clearly no.



The aggregate results demonstrate substantial differences between datasets.



On Adult, CatBoost is the strongest model across most aggregate metrics. TabPFN remains competitive and performs particularly strongly relative to XGBoost and LightGBM, but it does not achieve the best overall results.



On Bank Marketing, TabPFN is the strongest model across Accuracy, Recall, F1-score, and ROC-AUC, with CatBoost leading only in Precision.



On Credit-G, CatBoost is again the strongest model across most aggregate metrics, while XGBoost achieves the highest Precision.



The statistical analysis reinforces this dataset dependence. Of the 45 paired TabPFN-versus-tree comparisons, 34 are statistically significant at the unadjusted 0.05 level and 31 remain significant after Holm correction. However, significant differences occur in both directions depending on the dataset and baseline.



For example, TabPFN significantly outperforms CatBoost on several Bank Marketing metrics, while it significantly underperforms CatBoost on several Adult metrics. On Credit-G, none of the five TabPFN-CatBoost comparisons is statistically significant.



These findings support H3: the relative performance of TabPFN varies across datasets and evaluation metrics.



This also demonstrates why reporting only an overall average across datasets would be insufficient. Such an aggregate result could hide substantial differences in model behaviour between individual datasets.



\subsection{Research Question 4: What Computational Trade-offs Exist?}



The computational results reveal a clear trade-off between predictive performance and execution time.



TabPFN has the highest mean training time across all three datasets in this benchmark. Its mean training time is approximately 0.679 seconds on Adult, 0.685 seconds on Bank Marketing, and 0.659 seconds on Credit-G.



The difference becomes substantially larger for prediction time.



On Adult, TabPFN requires approximately 9.036 seconds for prediction compared with less than 0.1 seconds for each of the three tree-based models. A similar pattern occurs on Bank Marketing, where TabPFN requires approximately 8.946 seconds compared with approximately 0.04–0.09 seconds for the tree models.



Credit-G shows the same direction, although the absolute difference is smaller: TabPFN requires approximately 0.667 seconds compared with approximately 0.01 seconds for the tree-based models.



Therefore, the predictive advantages observed for TabPFN, particularly on Bank Marketing, are accompanied by a substantial prediction-time cost in this implementation.



This finding is important from a practical perspective. If predictive performance is the primary objective and the additional computation is acceptable, TabPFN may be attractive for some tasks. However, applications requiring very low inference latency may favour conventional tree-based models even when their predictive performance is slightly lower.



The computational results therefore support H4: differences in predictive performance are accompanied by meaningful computational trade-offs.



\subsection{Interpretation of the Model Differences}



The results suggest that TabPFN and gradient-boosted tree models should not be viewed simply as competing implementations of the same learning strategy.



The tree-based models learn a task-specific ensemble from the available observations. Their performance can therefore benefit directly from additional task-specific training data.



TabPFN follows a different paradigm by applying a pretrained model to a new tabular task. Its strong performance in several low- and medium-data configurations is consistent with the potential value of information acquired during pretraining when relatively little task-specific data are available.



However, the results also show that pretraining does not eliminate dataset dependence. CatBoost remains highly competitive on Adult and Credit-G, while TabPFN performs particularly strongly on Bank Marketing.



The differences between datasets therefore suggest that the usefulness of a pretrained tabular model depends on more than sample size alone. Feature characteristics, class distributions, and the relationship between predictors and target may all influence the relative suitability of the approaches.



The present benchmark does not isolate these individual dataset characteristics experimentally, so causal conclusions about why one dataset favours a particular model cannot be established from the current results. Nevertheless, the observed variation provides evidence that dataset-specific behaviour is an important consideration when selecting between tabular foundation models and gradient-boosted trees.



\subsection{Implications for Model Selection}



The results have several practical implications.



First, TabPFN should not automatically replace established gradient-boosted tree models. CatBoost achieves the strongest aggregate performance on two of the three datasets, demonstrating that conventional methods remain highly competitive.



Second, TabPFN should not be dismissed simply because tree-based models are strong baselines. Bank Marketing demonstrates that TabPFN can provide meaningful improvements across several predictive metrics and training-data regimes.



Third, model selection should consider the objective of the application. If predictive discrimination is the primary concern and the additional computational cost is acceptable, TabPFN may be a strong candidate. If inference speed is important, the considerably lower prediction times of the tree-based models may be an important practical advantage.



Finally, the results suggest that evaluating a model at a single training size is insufficient for understanding its practical behaviour. Learning curves reveal changes in relative performance that cannot be observed from a single aggregate benchmark point.



\subsection{Evaluation of the Hypotheses}



\subsubsection{H1: TabPFN will be competitive or superior when training data are limited.}



\textbf{Partially supported.}



TabPFN performs strongly in the smaller training regimes and outperforms XGBoost and LightGBM on several low-data comparisons. However, CatBoost remains competitive or superior on some datasets, particularly Adult and Credit-G.



Therefore, the results support the competitiveness of TabPFN in low-data settings but do not establish universal superiority.



\subsubsection{H2: TabPFN's relative advantage will change as training size increases.}



\textbf{Partially supported.}



The Adult results show a change in relative performance, with CatBoost becoming stronger at larger training sizes. However, Bank Marketing does not follow this pattern because TabPFN maintains its advantage across the evaluated training regimes.



The appropriate conclusion is therefore that training-data size influences relative model performance, but the direction and magnitude of the effect depend on the dataset.



\subsubsection{H3: TabPFN's performance relative to tree models will vary across datasets and metrics.}



\textbf{Supported.}



This is the clearest finding of the study. Adult, Bank Marketing, and Credit-G produce substantially different performance patterns, and the best model varies across evaluation metrics.



The statistical comparisons further demonstrate that the direction of model differences changes according to the dataset and baseline.



\subsubsection{H4: Predictive advantages will involve computational trade-offs.}



\textbf{Supported.}



TabPFN provides strong predictive performance in several settings but has substantially higher measured prediction times than the tree-based models in this benchmark. The results therefore demonstrate a clear predictive-performance versus computational-cost trade-off.



\subsection{Limitations}



Several limitations should be considered when interpreting the findings.



First, the study evaluates only three datasets. Although Adult, Bank Marketing, and Credit-G provide variation in dataset size and classification characteristics, they cannot represent the full diversity of real-world tabular machine-learning problems.



Second, only one tabular foundation model is evaluated. The findings therefore concern TabPFN specifically and should not automatically be generalized to all tabular foundation models.



Third, the benchmark evaluates binary classification only. The conclusions may not transfer directly to multiclass classification, regression, ranking, or other tabular prediction problems.



Fourth, Credit-G has a substantially smaller training partition than Adult and Bank Marketing. As a result, it could not be evaluated across the same large-data training regimes. Comparisons between datasets should therefore be interpreted with this difference in experimental coverage in mind.



Fifth, the computational measurements are specific to the experimental environment and implementation used for the benchmark. Training and prediction times can depend on hardware, software versions, preprocessing, model configuration, and system load. The reported timings should therefore be interpreted as measurements of this experimental setup rather than universal computational costs.



Finally, the benchmark does not attempt to identify the underlying causal characteristics that make one dataset more favourable to TabPFN than another. The observed dataset-level differences demonstrate that such variation exists, but additional controlled experiments would be required to determine which feature or dataset properties are responsible.



\subsection{Overall Interpretation}



The central finding of this study is that there is no universal winner between TabPFN and gradient-boosted tree algorithms.



TabPFN demonstrates strong predictive capability and provides its clearest advantage on Bank Marketing, where it performs strongly across multiple metrics and training-data regimes. At the same time, CatBoost remains highly competitive and achieves the strongest aggregate results on Adult and Credit-G.



The results therefore support a complementary view of tabular foundation models. TabPFN represents a valuable addition to the set of available tabular-learning methods, particularly when its predictive advantages justify its computational requirements. However, the continued strength of gradient-boosted tree algorithms means that they remain important baselines and practical alternatives.



More broadly, the study demonstrates the importance of evaluating tabular models across multiple training-data regimes rather than relying on a single benchmark configuration. Model rankings can change with training size, and the effect is not necessarily consistent across datasets.



The evidence consequently supports a dataset- and application-dependent approach to model selection. Rather than asking whether TabPFN is universally better than gradient-boosted trees, a more useful practical question is under which conditions the additional capabilities of a tabular foundation model provide sufficient predictive benefit to justify its computational cost.
